% ---------------------------------------------------------------------------------------
% MODELO DE TCC - INSTITUTO FEDERAL DE SERGIPE (IFS)
% ---------------------------------------------------------------------------------------
% Autor do Modelo: Mikael dos Santos Gois (mikaelgois1@gmail.com)
% Baseado nos modelos de: 
%   - Diogo D. Moreira (TCC, Dissertação e Tese)
%   - DECOM/UFOP (Monografia de Ciência da Computação)
% Modelo Base: abnTeX2 adaptado para as normativas internas de TCC do IFS (Art. 31 e 34)
% Versão: 1.0 (2026)
% ---------------------------------------------------------------------------------------
% Este arquivo faz parte de um modelo opensource para estudantes do IFS.
% Não remova os créditos dos autores originais e do modelo base.
% Este modelo é fornecido "no estado em que se encontra", sem garantias de qualquer tipo.
% ---------------------------------------------------------------------------------------

% [O que aparece no PDF Bookmark] {O que aparece impresso no documento}
\chapter[Introdução]{INTRODUÇÃO} \label{sec:introducao}
A introdução é a seção inicial do trabalho, onde o tema é apresentado e contextualizado.
Nela, o autor deve expor a motivação para a escolha do tema, os objetivos do estudo, a justificativa da relevância do trabalho e a metodologia adotada para alcançar os resultados esperados.
A introdução deve ser clara e objetiva, fornecendo ao leitor uma visão geral do que será abordado no restante do trabalho.
\cite{kitchenham2007guidelinesOnline}

\cite[text]{kitchenham2007guidelines}